\subsection{Extending Canonicalization Frames}\label{sec:canon-frame-work}

% TODO: compute cardinality bounds for canonicalization frames

% TODO: extend results on frame cardinalities and robust frames here
Following \cite{dym2024equivariant}, we make the following definitions in the context of canonicalization frames.

\begin{definition}
    Given a $G$-space $X$, a \emph{weighted canonicalization frame} is a function $\nu_\bullet : X \to \PP(X)$ where $\PP(X)$ is the space of Borel probability measures on $X$.
\end{definition}

\begin{definition}
    A weighted canonicalization frame $\nu_\bullet$ is \emph{orbital} if $\nu_x(X \sminus Gx) = 0$ for all $x \in X$.
\end{definition}

\begin{definition}
    A weighted canonicalization frame $\nu_\bullet$ is \emph{invariant} if $\nu_gx = \nu_x$ for all $x \in X$ and $g \in G$.
\end{definition}

% Going forward, we will assume that all weighted canonicalizations are orbital and invariant and thus not specify them as such.

\begin{proposition}
    Let $\nu_\bullet$ be an invariant weighted canonicalization frame.
    Then for any Borel function $f : X \to Y$, the function $\II_\nu (f)$ given by
    \[
        \II_\nu (f) (x) := \int_X f d \nu_x
    \]
    is $G$-invariant.
    \begin{proof}
        Let $g \in G$ and $x \in X$, and we simply compute
        \[
            \II_\nu (f) (gx)
            = \int_X f d \nu_{gx}
            = \int_X f d \nu_x
            = \II_\nu (f) (x).
        \]
    \end{proof}
\end{proposition}

\begin{definition}
    An invariant orbital weighted canonicalization frame $\nu_\bullet$ is continuous at $x$ if for any sequence $\seq{x} \sub X$ converging to $x$, we have $\nu_{x_n} \wto \nu_x$, i.e. for all $f : X \to \R$
    \[
        \int_X f d \nu_{x_n} \to \int_X f d \nu_x.
    \]
    A \emph{robust canonicalization frame} is a continuous invariant orbital weighted canonicalization frame.
\end{definition}

\begin{theorem}
    A weighted canonicalization frame $\nu_\bullet$ is continuous if and only if $\II_\nu$ preserves continuity.
    \begin{proof}
        First assume that $\nu_\bullet$ is continuous and let $x \in X$ with $\seq{x} \sub X$ such that $x_n \to x$.
        Then for any map $f : X \to \R$, we immediately have
        \[
            \II_\nu (f) (x_n)
            = \int_X f d \nu_{x_n}
            \to \int_X f d \nu_x
            = \II_\nu (f) (x)
        \]
        so we conclude that $\II_\nu$ preserves continuity.
        
        Now instead assume that $\II_\nu$ preserves continuity, let $f : X \to \R$ be a Borel function, and let $x \in X$ with $\seq{x} \sub X$ such that $x_n \to x$.
        Then by continuity of $\II_\nu (f)$ we have
        \[
            \int_X f d \nu_{x_n}
            = \II_\nu (f) (x_n)
            \to \II_\nu (f) (x)
            = \int_X f d \nu_x
        \]
        so we are done.
    \end{proof}
\end{theorem}


Similar to theorem \ref{thm:frame-equivalence}, we have a correspondence between weighted frames and weighted canonicalization frames.

\begin{theorem}\label{thm:weight-frame-equivalence}
    Let $X$ be a Polish space with a continuous left $G$-action.
    Then every weighted frame induces a weighted canonicalization frame and vice versa (though not necessarily via a bijection).
    In particular, this correspondence sends weakly equivariant weighted frames to invariant orbital weighted canonicalization frames and weighted canonicalization frames to equivariant weighted frames, and both directions preserve continuity.
    \begin{proof}
        To begin, let $\mu_\bullet$ be a weakly equivariant weighted frame.
        We define a weighted canonicalization frame $\nu_\bullet$ for Borel sets $E \sub X$ by
        \[
            \nu_x (E)
            = \Bar{\mu}_x(\set{g \in G : \inv{g}x \in E}).
        \]
        Then $\nu_\bullet$ is orbital and invariant since
        \[
            \nu_x (X \sminus Gx)
            = \Bar{\mu}_x(\set{g \in G : \inv{g}x \not\in Gx})
            = \Bar{\mu}_x(\emp) = 0.
        \]
        and
        \begin{align*}
            \nu_{hx} (E)
            &= \Bar{\mu}_{hx}(\set{g \in G : \inv{g}hx \in E}) \\
            &= \Bar{\mu}_x(\inv{h}\set{g \in G : \inv{g}hx \in E}) \\
            &= \Bar{\mu}_x(\set{g \in G : \pinv{hg}hx = \inv{g}x \in E}) \\
            &= \nu_x (E).
        \end{align*}
        Now assume $\mu_\bullet$ is continuous, and we wish to show that $\II_\nu (f)$ is continuous for any continuous $f$.
        To see this, first let $E \sub X$ be a Borel set and observe
        \[
            \int_X \chi_E d \nu_x
            = \nu_x (E)
            = \Bar{\mu}_x (\set{g \in G : \inv{g}x \in E})
            = \int_G \chi_E (\inv{g} x) d \Bar{\mu}_x (g).
        \]
        Since any Borel function $f : X \to \R$ can be approximated by simple functions, the Lebesgue Dominated Convergence Theorem gives us that 
        \[
            \int_X f d\nu_x = \int_G f(\inv{g} x) d \Bar{\mu}_x (g).
        \]
        But recall by lemma \ref{lem:weighted-frame-integral-equivalence} that 
        \[
            \int_G f(\inv{g} x) d \Bar{\mu}_x (g)
            = \int_G f(\inv{g} x) d \mu_x (g)
        \]
        so we have that $\II_\nu (f) = \II_\mu (f)$ and hence $\nu_\bullet$ is continuous.
        
        For the other direction, let $\nu_\bullet$ be an invariant orbital weighted canonicalization frame.
        We define a weighted frame $\mu_\bullet$ as follows.
        \footnote{The idea for this construction was suggested in part by Claude via \href{https://claude.ai/share/d975e47f-d488-4821-8653-ae7bbac49887}{this chat}.}
        First, fix some $x \in X$, and define a linear map $T_x : \cont{G}{\R} \to \cont{G/G_x}{\R}$ on continuous functions by
        \[
            T_x (f) (G_x g) := \int_{G_x} f(hg) dh
        \]
        where the integral is with respect to the Haar measure on $G_x$.
        Observe that this is well-defined since for $g' = h'g$ with $h' \in G_x$ we have
        \[
            T_x (f) (G_x g')
            = \int_{G_x} f(hg') dh
            = \int_{G_x} f(hh'g) dh
            = \int_{G_x} f(hg) dh
            = T_x (f) (G_x g)
        \]
        by the left-invariance of the Haar measure, and the integral is well-defined since $G$ is compact (hence all continuous functions are integrable).
        Additionally, $T_x$ is clearly linear by linearity of integration.
        Next, define a linear functional $L_x : \cont{G}{\R} \to \R$ by
        \[
            L_x (f) := \int_{Gx} T_x (f \circ q) d \nu_x
        \]
        where $q : Gx \to G/G_x$ is the quotient map.
        Observe that $L_x$ is indeed linear since $(f + g) \circ q = f \circ q + g \circ q$, and we have already established that $T_x$ is linear.
        Then by the Riesz-Markov-Kakutani Representation Theorem, there exists a unique Borel measure $\mu_x$ such that
        \[
            \int_G f d \mu_x = L_x (f) = \int_{Gx} T_x(f \circ q) d \nu_x,
        \]
        so we take all such $\mu_x$'s to be our $\mu_\bullet$.
        To show that $\mu_\bullet$ is equivariant, let $g \in G$ and $x \in X$.
        Then observe
        \begin{align*}
            T_{gx} (f) (G_x g')
            &= \int_{G_{gx}} f(h g') dh \\
            &= \int_{g G_x \inv{g}} f(h g') dh \\
            &= \int_{G_x} f(g h \inv{g}  g') dh \\
            &= \int_{G_x} (f(g h g'))  dh \\
            &= T_x (f(g \cdot -))(G_x g')
        \end{align*}
        so we have
        \begin{align*}
            \int_G f d \mu_{gx}
            &= \int_{Ggx} T_{gx} (f \circ q) d \nu_{gx} \\
            &= \int_{Gx} T_{gx} (f(q(u))) d \nu_x (u) \\
            &= \int_{Gx} T_x(f(g \cdot q(u))) d \nu_x(u) \\
            % FILL IN HERE
            &= \int_G f (g \cdot u) d \mu_x (u) \\
            &= \int_G f d g_* \mu_x
        \end{align*}
    \end{proof}
\end{theorem}

% TODO: consider free G-action case

% TODO: construct (weighted) canonicalization frames for specific groups/spaces